% GENERATED FILE. Edit canonical lessons, then run npm run generate:books.
% canonical-source-hash: fnv1a64:5f2f17a25ca2b19e
% canonical-lessons: LA-C438-conscendo, LA-C438-vectura, LA-C438-nauclerus, LA-C438-velificor, LA-C438-portitor, LA-R438-review-431-438

\chapter{To go on board, Transport, Ship's master, To set sail, Ferryman}
\label{ch:la-to-go-on-board-transport-ship-s-master-to-set-sail-ferryman}
\hlchaptermodality{438}

\begin{chapteropening}
\textbf{By the end of this chapter:} \emph{I can talk about taking a ship: going on board, transport and its cost, the ship's master, setting sail, and the ferryman.}

Complete the last lesson of chapter 438: I can talk about taking a ship: going on board, transport and its cost, the ship's master, setting sail, and the ferryman.
\end{chapteropening}

\section[cōnscendō, cōnscendere]{cōnscendō, cōnscendere — to go on board, to embark; to climb}

\label{lesson:LA-C438-conscendo}



\begin{quote}
Before the new one: say the Latin for a rudder, then the Latin for to row.
\end{quote}

\subsection*{You'll want to know: cōnscendō, cōnscendere}
\textbf{cōnscendō, cōnscendere} — \textquotedblleft{}to go on board, to embark; to climb\textquotedblright{}.

\begin{quote}
Prīmā lūce mercātōrēs nāvem cōnscendunt.
\end{quote}

\begin{quote}
\emph{At first light the merchants go on board the ship.}
\end{quote}

\emph{cōnscendō, cōnscendere, cōnscendī, cōnscēnsum}. \emph{Nāvem cōnscendere}, to embark; \emph{in nāvem} is also found.

\begin{cousinweb}
\emph{cōnscendō} is built from \emph{con-} and \emph{scandō} 'to climb' (EtymDB; Lewis \& Short): to climb up into a ship.
\end{cousinweb}

\begin{grammarlens}[title={What you've built}]
One more word for telling, explaining and describing at B1.
\end{grammarlens}

\subsection*{Your turn}
\begin{itemize}
\raggedright
  \item \emph{Say it:} \emph{cōnscendō, cōnscendere}
  \item \emph{Say it:} \emph{cōnscendō, cōnscendere}, once more
  \item \emph{Say it:} a sentence of your own with \emph{cōnscendō, cōnscendere}
  \item \emph{Recall:} say the Latin for a rudder, then the Latin for to row, then say \emph{cōnscendō, cōnscendere} again
\end{itemize}

\begin{tcolorbox}[breakable,colback=teal!4,colframe=teal!35!black,title={Before you move on}]
What does cōnscendō, cōnscendere mean? (\textquotedblleft{}To go on board, to embark; to climb\textquotedblright{}.) Say it once more.
\end{tcolorbox}
\section[vectūra, vectūrae]{vectūra, vectūrae — transport, carriage (of goods or people); the fare}

\label{lesson:LA-C438-vectura}



\begin{quote}
Before the new one: say the Latin for to row, then the Latin for to go on board.
\end{quote}

\subsection*{You'll want to know: vectūra, vectūrae}
\textbf{vectūra, vectūrae} — \textquotedblleft{}transport, carriage (of goods or people); the fare\textquotedblright{}.

\begin{quote}
Mercātor prō vectūrā mercium nummōs dat.
\end{quote}

\begin{quote}
\emph{The merchant pays coins for the transport of his goods.}
\end{quote}

First declension, feminine: \emph{vectūra, vectūrae f.}, from \emph{vehō} 'to carry'. Also the price paid for it.

\begin{cousinweb}
\emph{vectūra} is built on \emph{vehō} 'to carry' (Lewis \& Short), from its participle \emph{vectus}, with the ending \emph{-ūra}.
\end{cousinweb}

\begin{grammarlens}[title={What you've built}]
One more word for telling, explaining and describing at B1.
\end{grammarlens}

\subsection*{Your turn}
\begin{itemize}
\raggedright
  \item \emph{Say it:} \emph{vectūra, vectūrae}
  \item \emph{Say it:} \emph{vectūra, vectūrae}, once more
  \item \emph{Say it:} a sentence of your own with \emph{vectūra, vectūrae}
  \item \emph{Recall:} say the Latin for to row, then the Latin for to go on board, then say \emph{vectūra, vectūrae} again
\end{itemize}

\begin{tcolorbox}[breakable,colback=teal!4,colframe=teal!35!black,title={Before you move on}]
What does vectūra, vectūrae mean? (\textquotedblleft{}Transport, carriage (of goods or people); the fare\textquotedblright{}.) Say it once more.
\end{tcolorbox}
\section[nauclērus, nauclērī]{nauclērus, nauclērī — a ship's master, a shipowner}

\label{lesson:LA-C438-nauclerus}



\begin{quote}
Before the new one: say the Latin for to go on board, then the Latin for transport.
\end{quote}

\subsection*{You'll want to know: nauclērus, nauclērī}
\textbf{nauclērus, nauclērī} — \textquotedblleft{}a ship's master, a shipowner\textquotedblright{}.

\begin{quote}
Nauclērus nautās vocat et nāvem parat.
\end{quote}

\begin{quote}
\emph{The ship's master calls the sailors and gets the ship ready.}
\end{quote}

Second declension, masculine: \emph{nauclērus, nauclērī m.} The man who owns or commands a merchant ship.

\begin{cousinweb}
\emph{nauclērus} was borrowed from Greek \emph{nauklēros} 'shipowner' (EtymDB; Lewis \& Short), built from \emph{naus} 'ship' and \emph{klēros} 'lot, portion'.
\end{cousinweb}

\begin{grammarlens}[title={What you've built}]
One more word for telling, explaining and describing at B1.
\end{grammarlens}

\subsection*{Your turn}
\begin{itemize}
\raggedright
  \item \emph{Say it:} \emph{nauclērus, nauclērī}
  \item \emph{Say it:} \emph{nauclērus, nauclērī}, once more
  \item \emph{Say it:} a sentence of your own with \emph{nauclērus, nauclērī}
  \item \emph{Recall:} say the Latin for to go on board, then the Latin for transport, then say \emph{nauclērus, nauclērī} again
\end{itemize}

\begin{tcolorbox}[breakable,colback=teal!4,colframe=teal!35!black,title={Before you move on}]
What does nauclērus, nauclērī mean? (\textquotedblleft{}A ship's master, a shipowner\textquotedblright{}.) Say it once more.
\end{tcolorbox}
\section[vēlificor, vēlificārī]{vēlificor, vēlificārī — to set sail, to sail}

\label{lesson:LA-C438-velificor}



\begin{quote}
Before the new one: say the Latin for transport, then the Latin for a ship's master.
\end{quote}

\subsection*{You'll want to know: vēlificor, vēlificārī}
\textbf{vēlificor, vēlificārī} — \textquotedblleft{}to set sail, to sail\textquotedblright{}.

\begin{quote}
Ventus secundus est, et nautae laetī vēlificantur.
\end{quote}

\begin{quote}
\emph{The wind is favourable, and the sailors set sail happily.}
\end{quote}

Deponent: \emph{vēlificor, vēlificārī, vēlificātus sum}. Also an active form \emph{vēlificō}.

\begin{cousinweb}
\emph{vēlificor} is built on \emph{vēlum} 'sail' and \emph{faciō} 'to make' (Lewis \& Short): to make sail. EtymDB lists it without a parent.
\end{cousinweb}

\begin{grammarlens}[title={What you've built}]
One more word for telling, explaining and describing at B1.
\end{grammarlens}

\subsection*{Your turn}
\begin{itemize}
\raggedright
  \item \emph{Say it:} \emph{vēlificor, vēlificārī}
  \item \emph{Say it:} \emph{vēlificor, vēlificārī}, once more
  \item \emph{Say it:} a sentence of your own with \emph{vēlificor, vēlificārī}
  \item \emph{Recall:} say the Latin for transport, then the Latin for a ship's master, then say \emph{vēlificor, vēlificārī} again
\end{itemize}

\begin{tcolorbox}[breakable,colback=teal!4,colframe=teal!35!black,title={Before you move on}]
What does vēlificor, vēlificārī mean? (\textquotedblleft{}To set sail, to sail\textquotedblright{}.) Say it once more.
\end{tcolorbox}
\section[portitor, portitōris]{portitor, portitōris — a ferryman, a boatman; a customs officer}

\label{lesson:LA-C438-portitor}



\begin{quote}
Before the new one: say the Latin for a ship's master, then the Latin for to set sail.
\end{quote}

\subsection*{You'll want to know: portitor, portitōris}
\textbf{portitor, portitōris} — \textquotedblleft{}a ferryman, a boatman; a customs officer\textquotedblright{}.

\begin{quote}
Portitor viātōrēs ad alteram rīpam vehit.
\end{quote}

\begin{quote}
\emph{The ferryman carries the travellers to the other bank.}
\end{quote}

Third declension, masculine: \emph{portitor, portitōris m.} Two words in one spelling: a ferryman, and the officer who takes harbour dues.

\begin{cousinweb}
Lewis \& Short list two words \emph{portitor}: the customs officer, built on \emph{portus} 'harbour', and the carrier or ferryman, from the root of \emph{portō} 'to carry'. EtymDB does not list it.
\end{cousinweb}

\begin{grammarlens}[title={What you've built}]
One more word for telling, explaining and describing at B1.
\end{grammarlens}

\subsection*{Your turn}
\begin{itemize}
\raggedright
  \item \emph{Say it:} \emph{portitor, portitōris}
  \item \emph{Say it:} \emph{portitor, portitōris}, once more
  \item \emph{Say it:} a sentence of your own with \emph{portitor, portitōris}
  \item \emph{Recall:} say the Latin for a ship's master, then the Latin for to set sail, then say \emph{portitor, portitōris} again
\end{itemize}

\begin{tcolorbox}[breakable,colback=teal!4,colframe=teal!35!black,title={Before you move on}]
What does portitor, portitōris mean? (\textquotedblleft{}A ferryman, a boatman; a customs officer\textquotedblright{}.) Say it once more.
\end{tcolorbox}
\section[(dialogue)]{Words from chapters 431–438 — a review}

\label{lesson:LA-R438-review-431-438}



\begin{quote}
Say the Latin for a rein and a leather strap.
\end{quote}

\subsection*{Your turn}
Say these aloud:

\begin{itemize}
\raggedright
  \item a rein, a leather strap, a bridle, to bridle, a whip
  \item to ride, a saddle cloth, a spur, to jump down, to kick
  \item a beast of burden, a nag, a she-mule, a pack saddle, to load
  \item a stable, a manger, hay, to neigh, a muleteer
  \item a vessel, a little boat, a skiff, a punt, a raft
  \item the stern, the prow, the keel, a rowing bench, for cargo
  \item a ship's rope, a rope, a ship's tackle, a rudder, to row
  \item to go on board, transport, a ship's master, to set sail, a ferryman
\end{itemize}

\begin{tcolorbox}[breakable,colback=teal!4,colframe=teal!35!black,title={Before you move on}]
A rein? (\textbf{habēna, habēnae}.) A leather strap? (\textbf{lōrum, lōrī}.) A bridle? (\textbf{frēnum, frēnī}.) To bridle? (\textbf{frēnō, frēnāre}.) A whip? (\textbf{scutica, scuticae}.) To ride? (\textbf{equitō, equitāre}.) A saddle cloth? (\textbf{ephippium, ephippiī}.) A spur? (\textbf{calcar, calcāris}.) To jump down? (\textbf{dēsiliō, dēsilīre}.) To kick? (\textbf{calcitrō, calcitrāre}.) A beast of burden? (\textbf{iūmentum, iūmentī}.) A nag? (\textbf{caballus, caballī}.) A she-mule? (\textbf{mūla, mūlae}.) A pack saddle? (\textbf{clītellae, clītellārum}.) To load? (\textbf{onerō, onerāre}.) A stable? (\textbf{stabulum, stabulī}.) A manger? (\textbf{praesēpe, praesēpis}.) Hay? (\textbf{faenum, faenī}.) To neigh? (\textbf{hinniō, hinnīre}.) A muleteer? (\textbf{mūliō, mūliōnis}.) A vessel? (\textbf{nāvigium, nāvigiī}.) A little boat? (\textbf{nāvicula, nāviculae}.) A skiff? (\textbf{scapha, scaphae}.) A punt? (\textbf{linter, lintris}.) A raft? (\textbf{ratis, ratis}.) The stern? (\textbf{puppis, puppis}.) The prow? (\textbf{prōra, prōrae}.) The keel? (\textbf{carīna, carīnae}.) A rowing bench? (\textbf{trānstrum, trānstrī}.) For cargo? (\textbf{onerārius, onerāria, onerārium}.) A ship's rope? (\textbf{rudēns, rudentis}.) A rope? (\textbf{fūnis, fūnis}.) A ship's tackle? (\textbf{armāmenta, armāmentōrum}.) A rudder? (\textbf{gubernāculum, gubernāculī}.) To row? (\textbf{rēmigō, rēmigāre}.) To go on board? (\textbf{cōnscendō, cōnscendere}.) Transport? (\textbf{vectūra, vectūrae}.) A ship's master? (\textbf{nauclērus, nauclērī}.) To set sail? (\textbf{vēlificor, vēlificārī}.) A ferryman? (\textbf{portitor, portitōris}.)
\end{tcolorbox}
